% !TEX root = ../main.tex
\section{Conclusiones}

Se aplicó un flujo completo de minería de datos sobre el dataset Palmer Penguins comparando tres algoritmos de clustering no supervisado. La hipótesis inicial --- que algoritmos sensibles a la geometría global recuperarían mejor la estructura biológica que algoritmos basados en densidad --- queda confirmada por las métricas: Agglomerative Clustering con linkage de Ward obtiene un ARI de 0.916 y, tras mapear cada cluster a su especie mayoritaria (Cluster~0~$\to$~\emph{Adelie}, Cluster~1~$\to$~\emph{Gentoo}, Cluster~2~$\to$~\emph{Chinstrap}), una exactitud del 96.8\,\% contra la verdad de campo; K-Means logra ARI 0.793 y DBSCAN, 0.584.

El dataset resultó ideal para el ejercicio docente: tamaño manejable, tres especies bien definidas, presencia controlada de valores faltantes y suficiente solapamiento entre \emph{Adelie} y \emph{Chinstrap} como para que los algoritmos tengan que esforzarse. La estandarización previa de las cuatro variables fue clave; sin ella, \texttt{body\_mass\_g} (en gramos) dominaría las distancias frente a \texttt{bill\_depth\_mm} (en milímetros).

\subsection{Aportes del proyecto}

\begin{itemize}
  \item Notebook reproducible \texttt{notebook.ipynb} ejecutable en Jupyter local y en Google Colab sin pasos manuales.
  \item Comparativa cuantitativa de tres algoritmos con dos métricas complementarias --- interna y externa --- sobre el mismo pipeline.
  \item Cinco figuras vectoriales generadas por el propio notebook: pairplot, matriz de correlación, método del codo, proyección PCA y matriz de confusión.
  \item Presentación PowerPoint (5 diapositivas según rúbrica) generada programáticamente y sincronizada con las métricas del notebook.
  \item Repositorio organizado siguiendo Ciencia 2.0 (README, requirements, datos versionados, semillas fijadas).
\end{itemize}

\subsection{Limitaciones}

\begin{itemize}
  \item Sólo se usaron las cuatro variables numéricas. Las variables categóricas \texttt{sex} e \texttt{island} contienen información discriminante (por ejemplo, las islas de Gentoo y Chinstrap son disjuntas) y podrían mejorar la separación tras encoding.
  \item Los hiperparámetros de DBSCAN se fijaron a un valor estándar; una búsqueda fina con la curva $k$-distance probablemente lo equipararía a K-Means en este dataset.
  \item No se midió la estabilidad de los clusters frente a remuestreo. Un análisis con \emph{bootstrap} reforzaría la validez de los resultados \citep{hennig2007cluster}.
\end{itemize}

\subsection{Trabajo futuro}

Como continuación natural se contemplan: integración de las variables categóricas mediante \emph{one-hot encoding} o \emph{Gower distance}; comparación con Gaussian Mixture Models \citep{mclachlan2000gmm} para clusters elípticos con asignación probabilística; visualización no lineal con UMAP o t-SNE; y validación de estabilidad con \emph{bootstrap} y el índice de Jaccard \citep{hennig2007cluster}. En un horizonte más amplio, sería interesante replicar la metodología sobre el dataset \emph{Iris} y sobre \emph{Pima Indians Diabetes} para contrastar dominios.

\section*{Checklist de Ciencia 2.0}
\addcontentsline{toc}{section}{Checklist de Ciencia 2.0}

\begin{itemize}
  \item Dataset con licencia abierta (CC-0) y enlace al repositorio oficial.
  \item Código fuente íntegro en el notebook, sin dependencias propietarias.
  \item Semillas aleatorias fijadas (\texttt{RANDOM\_STATE=42}).
  \item \texttt{requirements.txt} con las librerías necesarias y sus versiones mínimas.
  \item Notebook ejecutable de principio a fin sin intervención manual.
  \item Reporte en \LaTeX{} con figuras vectoriales y citas verificables.
  \item Presentación de cinco diapositivas alineada con la rúbrica.
\end{itemize}
